\section{LEVEL 6 · Eindmissie}

\begin{missie}
Missie: voer een volledig onderzoek uit, los problemen op en leg je redenering uit. (LPD 1 – 4)
\end{missie}

\begin{doel}
\textbf{DOEL} \quad Na dit level kan je:

☐ een fenomeen uit de realiteit \textbf{beschrijven} met statistische begrippen (LPD 1);

☐ een probleem \textbf{aanpakken} met een heuristiek, wiskundig vertalen en terugvertalen met een antwoordzin (LPD 2);

☐ \textbf{ICT} gebruiken voor berekeningen en grafieken (LPD 3);

☐ je \textbf{redenering uitleggen} en fouten in redeneringen verbeteren (LPD 4).
\end{doel}

\begin{theorie}
\textbf{Aanpak van een probleem}

\begin{enumerate}
\item Noteer \textbf{gegeven en gevraagd}.
\item Maak een \textbf{schets of tabel}.
\item Deel op in \textbf{deelproblemen}, of onderzoek \textbf{alle mogelijkheden} (bijvoorbeeld het beste en het slechtste geval).
\item Vertaal wiskundig, reken uit en \textbf{vertaal terug} met een antwoordzin.
\item Controleer of je antwoord \textbf{realistisch} is.
\end{enumerate}
\end{theorie}

\subsection*{Eindmissie 1 · Mini-onderzoek: schermtijd en slaap}

Een leerlingenraad onderzoekt of schermtijd en slaap samenhangen. Zestien leerlingen van het zesde jaar, die zich vrijwillig aanmeldden, noteerden hun gemiddelde schermtijd en slaapduur (in uur per dag):

\breedtes{14/5/5/5/5/5/5/5/5/5/5/5/5/5/5/5/5}
\begin{tabular}{lcccccccccccccccc}
\hline
leerling & 1 & 2 & 3 & 4 & 5 & 6 & 7 & 8 & 9 & 10 & 11 & 12 & 13 & 14 & 15 & 16 \\
\hline
scherm & 1,5 & 2 & 2,5 & 3 & 3 & 3,5 & 4 & 4 & 4,5 & 5 & 5 & 5,5 & 6 & 6,5 & 7 & 7,5 \\
slaap & 8,2 & 8,5 & 7,8 & 8,0 & 7,6 & 7,9 & 7,2 & 7,7 & 7,0 & 7,3 & 6,6 & 7,1 & 6,4 & 6,9 & 6,0 & 6,3 \\
\hline
\end{tabular}

\textbf{a)} \emph{(LPD 1)} Formuleer een probleemstelling, een onderzoeksvraag en een hypothese. \lijn

\textbf{b)} Noteer de twee variabelen en het soort variabele en data. \lijn

\textbf{c)} \emph{(LPD 11)} De 16 leerlingen meldden zich vrijwillig aan. Is de steekproef representatief voor alle leerlingen van de school? Verklaar met twee begrippen uit deze bundel. \lijn

\textbf{d)} \emph{(LPD 3)} Bereken met ICT het gemiddelde en de mediaan van beide variabelen. Maak met Excel of GeoGebra een puntenwolk (spreidingsdiagram) met schermtijd op de $x$-as en slaap op de $y$-as. Welk verband zie je? \lijn

\textbf{e)} \emph{(LPD 12)} Mag je besluiten dat meer schermtijd voor minder slaap zorgt? Verklaar met een omgekeerde richting en een derde variabele. Hoe zou je causaliteit kunnen onderzoeken? \lijn

\textbf{f)} \emph{(LPD 4)} Schrijf een besluit van vier à vijf zinnen voor de directie: resultaat, beperkingen en advies. \lijn

\begin{oplossing}
\textbf{Oplossing Eindmissie 1} \quad a) (Open.) Bv. probleemstelling: sommige leerlingen zijn ’s ochtends erg moe. Onderzoeksvraag: ‘Hangt de schermtijd samen met de slaapduur?’ Hypothese: ‘Meer schermtijd gaat samen met minder slaap.’ b) Schermtijd (uur per dag) en slaapduur (uur per dag): kwantitatief; numerieke data. c) Nee: vrijwillige respons (alleen wie zich aanmeldt, mogelijk uitgesproken of geïnteresseerd) en een kleine steekproefomvang (16 leerlingen); mogelijk ook niet-overeenkomende samenstelling (enkel zesde jaar). d) Schermtijd: $\overline{x}\approx 4{,}4$ uur, Med = 4,25 uur; slaap: $\overline{x}\approx 7{,}3$ uur, Med = 7,25 uur. Puntenwolk: dalend, duidelijk negatieve samenhang (hoe meer schermtijd, hoe minder slaap). e) Nee: observationeel onderzoek toont enkel samenhang. Omgekeerd: wie moeilijk slaapt, zit langer op het scherm. Derde variabele: stress of veel huiswerk of een onregelmatig ritme beïnvloedt beide. Causaliteit: experiment met door loting gevormde groepen (bv. ene groep schermvrije avonden, andere groep niet) en dezelfde meting van de slaap. f) (Open.) Let op correcte cijfers, vermelding van vrijwillige respons en kleine steekproef, samenhang geen causaliteit, concreet advies (bv. een groter, aselect onderzoek).
\end{oplossing}

\subsection*{Eindmissie 2 · Twee problemen}

\subsubsection*{Probleem A \quad ★★★ \quad De enquête van de kinderopvang}

Een kinderopvang mailt een enquête naar 450 ouders. Het responspercentage is 64 \%. Van de ouders die antwoorden, vindt 75 \% de opvang betaalbaar.

\textbf{a)} Hoeveel ouders antwoordden? Hoeveel van hen vinden de opvang betaalbaar? \lijn

\textbf{b)} Mag de opvang besluiten dat 75 \% van alle 450 ouders de opvang betaalbaar vindt? Onderzoek het beste en het slechtste geval en schrijf een antwoordzin. \lijn

\begin{oplossing}
\textbf{Oplossing Probleem A} \quad a) $450\cdot 0{,}64=288$ ouders antwoordden; $288\cdot 0{,}75=216$ vinden het betaalbaar. b) Er zijn $450-288=162$ ouders die niet antwoordden. Slechtste geval: geen van hen vindt het betaalbaar: $\frac{216}{450}=48\,\%$. Beste geval: allemaal: $\frac{216+162}{450}=\frac{378}{450}=84\,\%$. Antwoordzin: door de non-respons ligt het aandeel in de hele groep ergens tussen 48 \% en 84 \%; ‘75 \% van alle ouders’ mag dus niet zomaar besloten worden.
\end{oplossing}

\subsubsection*{Probleem B \quad ★★★ \quad De prijsverhoging}

Een dag opvang stijgt van 12 euro naar 15 euro.

\textbf{a)} Met hoeveel procent stijgt de prijs? \lijn

\textbf{b)} Een politicus belooft: ‘We laten de prijs dalen met hetzelfde percentage.’ Welke prijs krijg je dan? Is dat opnieuw 12 euro? \lijn

\textbf{c)} Met hoeveel procent moet de prijs van 15 euro dalen om terug 12 euro te zijn? \lijn

\begin{oplossing}
\textbf{Oplossing Probleem B} \quad a) $\frac{15-12}{12}=0{,}25=25\,\%$. b) $15\cdot 0{,}75=11{,}25$ euro: niet 12 euro, want de daling wordt op het hogere bedrag (15) berekend. c) $\frac{12-15}{15}=-0{,}20=-20\,\%$: 20 \% daling.
\end{oplossing}

\subsection*{Eindmissie 3 · Een krantenbericht onder de loep}

\begin{voorbeeld}
\textbf{Gamers halen slechtere punten: spelcomputers zijn de schuldige!} Een online peiling bij 1\,250 bezoekers van de website van een speelgoedwinkel toont dat 85 \% van de ouders vindt dat gamen slecht is voor de schoolresultaten. Kinderen die meer dan drie uur per dag gamen, halen een gemiddeld rapportcijfer van 6,8 tegenover 7,6 voor niet-gamers. ‘Gamen veroorzaakt slechtere punten’, besluit de winkel, die houten speelgoed verkoopt.
\end{voorbeeld}

\textbf{a)} Noteer minstens drie problemen in dit bericht en benoem ze met de juiste begrippen. \lijn

\textbf{b)} Leg in je eigen woorden aan een ouder uit waarom de conclusie ‘gamen veroorzaakt slechtere punten’ niet bewezen is. \lijn

\textbf{c)} Hoe zou je onderzoeken of gamen de punten beïnvloedt? \lijn

\begin{oplossing}
\textbf{Oplossing Eindmissie 3} \quad a) Online peiling bij bezoekers van een webshop = vrijwillige respons en opportunistische steekproef (niet representatief); de opdrachtgever/winkel heeft belang bij het resultaat; samenhang wordt gelezen als causaliteit (derde variabele of omgekeerde richting). b) Dat gamers lagere punten halen, toont enkel samenhang. Misschien gamen kinderen met slechte punten meer, of is er een derde factor (weinig motivatie, weinig toezicht). Zonder experiment weet je niet wat oorzaak en gevolg is. c) Experiment met door loting gevormde groepen (een groep beperkt gamen, een controlegroep niet), dezelfde metingen van de punten en een voldoende grote, representatieve steekproef.
\end{oplossing}
